\section{序列建模的两条路线各自付什么代价}
\label{sec:14-Deep-Learning/05-Attention-and-Sequence-Models/06}

你要把一个转写模型放进一台边缘设备，硬性要求是：延迟恒定，内存恒定，会议开八小时也不能崩。你手上的 Transformer 在这条要求前直接出局——上一篇算过，缓存每个 token 半兆，随时长一路上涨。而 2015 年那套 LSTM 无论跑多久，隐状态永远是 512 个浮点数，每步开销一模一样。

这不是``旧方法在某个角落还有余温''。把两条路线摆在同一组坐标上看，它们各自在三个维度上一赢一输，且赢和输是同一个结构决定的。

\subsection{三个维度，两组符号相反的答案}

\emph{第一个维度是梯度路径长度}。RNN 的隐状态递推 \(h_t=f(h_{t-1},x_t)\) 让相距 \(k\) 步的两个位置之间只有一条通路，梯度必须逐步倒回去：

\[
\frac{\partial h_t}{\partial h_{t-k}} = \prod_{j=t-k+1}^{t} \frac{\partial h_j}{\partial h_{j-1}} .
\]

\(k\) 个雅可比连乘，量级由它们的谱半径 \(\lambda\) 决定，大致是 \(\lambda^{k}\)。\hyperref[sec:03-Linear-Algebra/05-Eigenvalues-and-Eigenvectors/02]{``对角化与动力系统：把耦合演化变成独立缩放''}讲的正是这种幂次行为：小于 1 的方向按几何级数塌陷，大于 1 的方向按几何级数爆炸，只有恰好等于 1 的方向能把信息原样传下去，而这是一个测度为零的巧合。\(k=200\) 时 \(0.95^{200}\approx 3\times10^{-5}\)——不是学得慢，是梯度到不了。LSTM 与 GRU 的门控本质上是给细胞状态开一条加法通路，把某些方向的 \(\lambda\) 人为钉在 1 附近，和 \hyperref[sec:14-Deep-Learning/05-Attention-and-Sequence-Models/04]{``残差与归一化让深层可训''}里的恒等项是同一个想法用在时间轴上。它缓解了问题，没有取消 \(\lambda^k\) 这个形式。

Transformer 这一维上的答案完全不同。位置 \(t\) 直接从位置 \(s\) 读取，中间不经过任何位置：

\[
\frac{\partial z_t}{\partial v_s} = A_{ts}\,I .
\]

一个标量权重，没有连乘。\(s\) 与 \(t\) 相距 1 还是相距 10000，这个式子的形式完全一样。路径长度是常数 1，与距离无关——这才是 Transformer 真正的结构优势，比``并行''更根本。

这里要说准确一点，否则会滑成一句口号。路径短不等于梯度大：\(A_{ts}\) 本身可能是 \(10^{-6}\)，信息照样传不动，只不过这时是模型学出来的选择，不是结构强加的衰减。而且 Transformer 并没有消灭连乘，它只是把连乘从\emph{距离方向}挪到了\emph{深度方向}——\(L\) 层仍然要连乘 \(L\) 个雅可比。区别在于 \(L\) 是你定的常数，\(k\) 是数据给的，可以是几万。把一个随数据增长的指数换成一个固定的指数，再用残差把后者摁住，这就是全部的技巧。

\emph{第二个维度是训练时的并行性}。\(h_t\) 要等 \(h_{t-1}\) 算完，长度 \(n\) 的序列在时间轴上就有 \(n\) 个严格串行的步骤，无论多少张卡都压不下去。Transformer 训练时所有位置的 \(Q,K,V\) 一次投影完成，注意力就是两次大矩阵乘法，时间维上的串行步数是常数。这一点常被说成``更快''，更准确的说法是：同样的墙钟时间里能走过更多数据。Transformer 在数据规模上的优势有相当一部分来自这里，而不是来自单步建模能力更强。

\emph{第三个维度是推理成本，符号反过来了}。RNN 的状态是固定维的，每一步只做一次 \(O(d^2)\) 的更新，与已经生成了多长完全无关，内存占用是常数。Transformer 每一步要读全部历史，计算 \(O(nd)\)，内存 \(O(nd)\) 且持续增长。开头那个八小时会议的场景，输的正是这一维。

\begin{center}
\begin{tikzpicture}[scale=0.95]
% 图例
\fill[black!30] (0,0.95) rectangle (0.45,1.25);
\node[anchor=west] at (0.55,1.1) {\footnotesize RNN};
\fill[black!60] (2.0,0.95) rectangle (2.45,1.25);
\node[anchor=west] at (2.55,1.1) {\footnotesize Transformer};
\node[anchor=west] at (5.0,1.1) {\footnotesize 条越长，代价越大};
% 三行
\foreach \y/\lab/\wr/\wt/\tr/\tt in {%
  0/{梯度路径长度}/5.4/0.55/{\(O(k)\)}/{\(O(1)\)},
  -1.35/{训练时的串行步数}/5.4/0.55/{\(O(n)\)}/{\(O(1)\)},
  -2.7/{每步推理成本}/0.55/4.7/{\(O(1)\)}/{\(O(n)\)}}{
  \node[anchor=east] at (-0.15,{\y-0.02}) {\footnotesize \lab};
  \fill[black!30] (0,{\y+0.08}) rectangle (\wr,{\y+0.42});
  \fill[black!60] (0,{\y-0.42}) rectangle (\wt,{\y-0.08});
  \node[anchor=west] at ({\wr+0.12},{\y+0.25}) {\footnotesize \tr};
  \node[anchor=west] at ({\wt+0.12},{\y-0.25}) {\footnotesize \tt};
}
\draw[black!70] (0,0.62) -- (0,-3.2);
\end{tikzpicture}
\end{center}

\emph{同样三个维度，两条路线的长短条正好互补：Transformer 在梯度路径和训练并行上完胜，在每步推理成本上完败。没有哪一行是``技术进步''的方向——三行合起来才是一次结构选择的全貌。}

图上的第三行解释了为什么这场比较没有终局。前两行赢下来的东西让 Transformer 吃得下互联网规模的数据，所以它赢了训练；第三行输掉的东西在部署侧一天比一天贵，因为窗口只会越拉越长。

\subsection{状态空间模型想同时拿到前两行和第三行}

如果一条路线能在三行上都取短条，那它就该赢。S4、Mamba 一类的状态空间模型就是奔着这个去的，它们的切入点是把递推限制成\emph{线性}的：

\[
h_t = A h_{t-1} + B x_t,\qquad y_t = C h_t .
\]

去掉非线性之后，递推可以整体展开：\(h_t=\sum_{j\le t}A^{t-j}Bx_j\)，也就是输入与核 \(\{CA^{i}B\}\) 的一次卷积。卷积可以用 \hyperref[sec:03-Linear-Algebra/08-Complex-Matrices-and-Fourier-Structure/03]{``Fourier 基与 FFT''}里的办法在 \(O(n\log n)\) 内一次算出全部输出，训练时的串行步数因此坍缩成常数——第二行拿下了。推理时换回递推形式，状态固定维、每步常数开销——第三行也拿下了。梯度方向上，\(A\) 的谱被专门设计（HiPPO 一类的初始化）使得长程分量不塌陷，第一行至少不吃亏。

三行都短，代价在哪？在一个这张图上没有的第四行：\emph{状态容量}。

无论历史多长，\(h_t\) 始终是 \(d_{\text{state}}\) 个数。把长度 \(n\) 的历史压进固定维的向量，\(n\) 一大就必然有损——这不是模型没训好，是容量本身的限制。KV cache 之所以昂贵，正因为它是\emph{无损}的：每个 token 的 Key 与 Value 原样保留，任何时刻都能被精确检索回来。上一篇算的那笔线性增长的显存账，买的就是这份无损。

这个区别在合成任务上可以直接测量。给模型一串随机键值对，然后问某个键对应什么值，需要的是精确回忆而不是概括。注意力模型可以做对任意长度，因为对应的 Value 就存在缓存里；固定状态的模型在序列长度超过状态容量之后准确率系统性下滑。这类实验反复出现，也是当前混合架构的直接动机——大部分层用 SSM 拿走效率，插进少数几层注意力保住精确检索。这是一条经验路线，不同任务上最优的混合比例并不一致。

\subsection{没有免费的记忆}

把六篇拉通看，这一章其实只讲了一次交易和三笔付款。

注意力用``内容决定的加权读取''换掉了固定的连接结构：读谁不再由位置写死，而由内容当场算出来。这是第一篇的收益。代价随后一笔一笔到账——位置信息在等变性下被抹平，必须外挂；恒等通路在深度上不能断，否则堆不上去；全局可读的标价是 \(n^2\) 的计算和线性增长的缓存。多头这一篇看似是收益，其实也是一次分配：同一笔参数预算切成几条低秩通道，秩上界就是分工的价钱。

最后这一篇把交易本身放回了选择空间。记忆有两种付法：按长度付内存，把历史无损地存着，代价随时间线性上涨；或者按压缩付信息，状态固定但历史必然失真。SSM 与注意力站在这条线的两端，混合架构在中间挑一个点。哪种更好取决于任务需要的是概括还是回忆，取决于部署环境愿意为哪一种代价买单。

一个模型能记住什么，最终由它愿意为记忆付多少钱决定。这个约束不来自哪个具体架构，短期内也不会有哪个新架构把它取消。
